\section{An entire normalization for all independent harmonic exponents}
\label{eht:sec}

The independent-inner-exponent question in the incoming exact-identities
report connects directly to an established interpolant of truncated
multiple zeta functions.  We identify that interpolant exactly, give a
self-contained residue proof of its entire continuation, and develop a
relative version with two endpoints.  The resulting endpoint, Gamma, and
Stieltjes identities give an analytic replacement for the one-parameter
Gamma normalization when the inner exponents move independently.

The result does not assert a reduction of every ordered coefficient to
ordinary zeta values.  In particular the nonsymmetric harmonic Stieltjes
coordinate from the incoming ordered-resonance article remains visible.
The analytic continuation and residue structure of multiple Hurwitz zeta
functions, and the quasi-shuffle algebra, are classical
\cite{eht:MehtaViswanadham,eht:Hoffman}.  The one-endpoint entire
interpolation and its stuffle law are due to Ihara, Nakamura, and Yamamoto,
with earlier overlapping interpolation work credited in their paper
\cite{eht:INY}.  The relative ordered quotient, residue calculation, and
explicit consequences below organize this known construction for the
repository's present questions.  No literature-wide priority claim is made.

\subsection{Definitions and the finite formula}

For $\Re a>0$ and a word $w=(s_1,\ldots,s_r)$, let
\begin{equation}\label{eht:Z}
 Z_a(w)=\sum_{n_1>\cdots>n_r\ge0}
              \prod_{j=1}^r(n_j+a)^{-s_j},\qquad Z_a(\varnothing)=1.
\end{equation}
Use principal powers.  The series is initially normally convergent when
$\Re(s_1+\cdots+s_j)>j$ for every $j$.  Its meromorphic continuation is
understood below.  The deliberately smaller domain $\Re s_j>1$ for every
$j$ suffices for all initial product manipulations.

For $\Re a,\Re b>0$, define recursively
\begin{align}
 \mathcal H_{a,b}(\varnothing)&=1,\notag\\
 \mathcal H_{a,b}(s_1,\ldots,s_r)
 &=Z_b(s_1,\ldots,s_r)
   -\sum_{j=1}^rZ_a(s_1,\ldots,s_j)
                    \mathcal H_{a,b}(s_{j+1},\ldots,s_r).
 \label{eht:rec}
\end{align}
The sum includes the final term $Z_a(s_1,\ldots,s_r)$.  Equivalently,
\begin{equation}\label{eht:finite}
 \mathcal H_{a,b}(w)=
 \sum_{q=0}^r(-1)^q
 \sum_{0=i_0<i_1<\cdots<i_q\le r}
 \left\{\prod_{\nu=1}^q
   Z_a(s_{i_{\nu-1}+1},\ldots,s_{i_\nu})\right\}
 Z_b(s_{i_q+1},\ldots,s_r).
\end{equation}
For $q=0$ the summand means $Z_b(w)$.  There are exactly $2^r$ terms,
counting repetitions before simplification.  Thus the definition requires
no limiting choice, ray, or finite-part prescription.

The first two cases are
\begin{align}
 h_{a,b}(s):=\mathcal H_{a,b}(s)
   &=\zeta(s,b)-\zeta(s,a),\label{eht:h}\\
 \mathcal H_{a,b}(s,t)
   &=Z_b(s,t)-Z_a(s,t)-\zeta(s,a)
                         \{\zeta(t,b)-\zeta(t,a)\}.
 \label{eht:depthtwo}
\end{align}
The second line already shows why subtracting just $Z_a-Z_b$ fails:
its residue at $s=1$ is generally nonzero.

There is also a shorter form in terms of the weakly ordered Hurwitz
function $Z_a^\star$, whose inequalities in \eqref{eht:Z} are all weak:
\begin{equation}\label{eht:star}
 \mathcal H_{a,b}(w)=\sum_{j=0}^r(-1)^j
       Z_a^\star(s_j,\ldots,s_1)Z_b(s_{j+1},\ldots,s_r).
\end{equation}
The empty prefix has value one.  The inverse of the strict ordered series
has coefficient $(-1)^jZ_a^\star(s_j,\ldots,s_1)$: expand each inverse
one-index factor, then reverse the weak order.  This proves
\eqref{eht:star} in an absolute-convergence domain and hence meromorphically.
Proposition 2.1 of \cite{eht:INY} therefore gives the exact identification
\begin{equation}\label{eht:INY}
                 \mathcal H_{a,1}(w)=\Psi(w;a-1).
\end{equation}
Their Corollary 4.2 already proves that this is entire in all indices;
the proof below provides a different route through explicit prefix
residues and treats the relative endpoints together.

\begin{theorem}[Entire relative harmonic interpolation]\label{eht:main}
For every depth $r$, the expression \eqref{eht:finite} extends to an
entire function of $(s_1,\ldots,s_r)\in\mathbb C^r$.  It is jointly
holomorphic with $a,b$ on the right half-plane.  Every fixed mixed spectral
derivative is consequently an ordinary derivative of one entire function,
including where several polar hyperplanes of its separate summands meet.

If $N\ge0$ is an integer and $a=b+N$, then
\begin{equation}\label{eht:interval}
 \boxed{\mathcal H_{b+N,b}(s_1,\ldots,s_r)=
  \sum_{N>n_1>\cdots>n_r\ge0}
              \prod_{j=1}^r(n_j+b)^{-s_j}.}
\end{equation}
In particular $\mathcal H_{M,1}$ is the strict finite multiple harmonic
sum with largest index at most $M-1$.
\end{theorem}

\subsection{A shift-independent prefix-residue lemma}

Write $S_j=s_1+\cdots+s_j$.  Introduce the rational-polynomial functions
\begin{equation}\label{eht:beta}
 \beta_0(z)=\frac1{z-1},\qquad
 \beta_1(z)=-\frac12,\qquad
 \beta_{2q}(z)=\frac{B_{2q}}{(2q)!}(z)_{2q-1}\quad(q\ge1),
 \qquad \beta_{2q+1}(z)=0\quad(q\ge1).
\end{equation}
Here the $B_{2q}$ are Bernoulli numbers, with $B_1=-1/2$.

\begin{lemma}[Prefix residue factorization]\label{eht:residue}
The possible polar divisors of $Z_a(s_1,\ldots,s_r)$ are among
\[
 s_1=1,\qquad S_k=k-m\quad(2\le k\le r,\ m\in\mathbb Z_{\ge0}).
\]
At a generic point of each divisor the pole is at most simple.  In the
normal coordinate $S_k-k+m$, its residue is
\begin{equation}\label{eht:resfactor}
 C_{k,m}(s_1,\ldots,s_k)Z_a(s_{k+1},\ldots,s_r),
\end{equation}
where $C_{1,0}=1$ and, for $k\ge2$,
\begin{equation}\label{eht:C}
 C_{k,m}=
 \sum_{\substack{e_1+\cdots+e_{k-1}=m\\e_i\ge0}}
 \prod_{i=1}^{k-1}
 \beta_{e_i}\left(S_i-(i-1)+\sum_{j<i}e_j\right),
 \qquad\text{restricted to }S_k=k-m.
\end{equation}
The coefficient is independent of $a$.  A vanishing $C_{k,m}$ simply means
that the indicated possible divisor is absent.
\end{lemma}

\begin{proof}
Sum first over $n_1$, put $x=n_2+a$, and use a finite
Euler--Maclaurin expansion
\begin{equation}\label{eht:EM}
 \zeta(z,x+1)=\frac{x^{1-z}}{z-1}-\frac12x^{-z}
    +\sum_{q=1}^Q\frac{B_{2q}}{(2q)!}(z)_{2q-1}x^{-z-2q+1}
    +R_Q(z,x).
\end{equation}
On any prescribed compact set of $z$, sufficiently large $Q$ makes the
remainder holomorphic there, with a uniform estimate
$R_Q(z,n+a)=O((n+1)^{1-\Re z-2Q})$ on compact right-half-plane sets of
$a$.  The same is true after fixed parameter derivatives, with logarithmic
factors.  This follows directly from the integral remainder with a bounded
periodic Bernoulli function.  In the remaining $r-1$ sums the number and
size of the inner terms have at most polynomial growth in $n_2$ on a
compact set of all exponents.  Taking $Q$ larger therefore makes the
remainder normally summable on that compact.

Each explicit term in \eqref{eht:EM} replaces the first two exponents by
$s_1+s_2-1+e$, with coefficient $\beta_e(s_1)$.  Repeat this step in the
lower-depth functions.  This proves meromorphic continuation by induction
and shows that every possible polar equation is a prefix equation of the
displayed form.  The pole at $s_1=1$ has residue
$Z_a(s_2,\ldots,s_r)$.

To reach a pole involving exactly the first $k$ indices, merge those
indices in $k-1$ steps.  Before step $i$ the current leading exponent is
$S_i-(i-1)+e_1+\cdots+e_{i-1}$.  After all merges, the final Hurwitz pole
has normal coordinate
$S_k-k+e_1+\cdots+e_{k-1}$.  The paths producing $S_k=k-m$ are exactly
the finite set in \eqref{eht:C}.  The remaining indices form the factor
$Z_a(s_{k+1},\ldots,s_r)$.  All coefficients in this procedure are
independent of the common shift.

The only denominators in the merging coefficients arise from proper
prefix hyperplanes.  Away from intersections, the final pole is simple.
Since $Q$ can be chosen arbitrarily large on a prescribed compact, the
holomorphic remainders cannot add further local poles.  This proves the
lemma, including its joint parameter assertion.
\end{proof}

For example,
\[
 C_{2,0}=\frac1{s_1-1},\qquad C_{2,1}=-\frac12,
 \qquad C_{2,2}=\frac{s_1}{12},\qquad C_{2,3}=0,
 \qquad C_{3,0}=\frac1{(s_1-1)(s_1+s_2-2)}.
\]
These are residues in a sum-of-indices coordinate; replacing that
coordinate by a scaled parameter rescales the residue.

\begin{proof}[Proof of Theorem~\ref{eht:main}]
Induct on the depth.  At depth one the common residue of the two Hurwitz
functions cancels.  Suppose all lower-depth $\mathcal H_{a,b}$ are entire.
Equation~\eqref{eht:rec} shows that the only remaining possible poles are
global prefix hyperplanes of the full word.  Fix one involving the first
$k$ indices and work away from all other hyperplanes.  The terms with
$j<k$ in \eqref{eht:rec} have no pole there.  By
Lemma~\ref{eht:residue}, the residue of the rest is
\[
 C_{k,m}\left\{
 Z_b(s_{k+1},\ldots,s_r)
 -\sum_{j=k}^rZ_a(s_{k+1},\ldots,s_j)
                  \mathcal H_{a,b}(s_{j+1},\ldots,s_r)
 \right\}=0.
\]
The bracket is exactly the lower-depth convolution identity
\eqref{eht:rec}, including the empty-prefix term at $j=k$.
Thus every generic polar hypersurface is removable.  The locally finite
union of their intersections has complex codimension at least two.
Hartogs extension removes these remaining intersections; equivalently,
one can cancel each linear factor in a local meromorphic denominator.
This proves joint entire continuation in the exponents.  The same proof
with $a,b$ in compact right-half-plane sets gives joint holomorphy.

For the interpolation assertion, take $\Re s_j>1$.  Split the strictly
ordered indices for $Z_b(w)$ at $N$: a uniquely determined initial block
lies in $n\ge N$, and the remaining block lies in $n<N$.  Hence
\[
 Z_b(w)=\sum_{j=0}^rZ_{b+N}(s_1,\ldots,s_j)
   \sum_{N>n_{j+1}>\cdots>n_r\ge0}
                    \prod_{\nu=j+1}^r(n_\nu+b)^{-s_\nu}.
\]
The triangular recursion identifies the finite sums with
$\mathcal H_{b+N,b}$.  Both sides are entire in all exponents, so the
identity extends to $\mathbb C^r$.
\end{proof}

\subsection{Composition, the stuffle product, and endpoint singularities}

For bookkeeping, attach a noncommuting letter $[s]$ to each exponent and
form the coefficientwise series
\[
 F_a=1+\sum_{r\ge1}\sum_{s_1,\ldots,s_r}
             Z_a(s_1,\ldots,s_r)[s_1]\cdots[s_r].
\]
This notation requires only a finite generating alphabet and its additive
closure at any application; it is not a sum over an uncountable set.
Concatenation of words gives
\begin{equation}\label{eht:quotient}
 \mathcal H_{a,b}=F_a^{-1}F_b.
\end{equation}

\begin{corollary}[Exact composition and shifts]\label{eht:transport}
With products interpreted by deconcatenation,
\begin{equation}\label{eht:cocycle}
 \boxed{\mathcal H_{a,b}(w)=
 \sum_{j=0}^r\mathcal H_{a,c}(s_1,\ldots,s_j)
                  \mathcal H_{c,b}(s_{j+1},\ldots,s_r).}
\end{equation}
Moreover
\begin{align}
 \mathcal H_{a+1,b}(w)
   &=\mathcal H_{a,b}(w)+a^{-s_1}
                              \mathcal H_{a,b}(s_2,\ldots,s_r),
       \label{eht:upper-shift}\\
 \mathcal H_{a,b}(w)
   &=\mathcal H_{a,b+1}(w)+b^{-s_r}
                              \mathcal H_{a,b+1}(s_1,\ldots,s_{r-1}).
       \label{eht:lower-shift}
\end{align}
All of these identities hold for arbitrary complex exponents.
\end{corollary}

\begin{proof}
The composition formula is
$F_a^{-1}F_cF_c^{-1}F_b=F_a^{-1}F_b$.
Splitting off the smallest index gives
$F_a=F_{a+1}E(a)$, where $E(a)=1+\sum_s a^{-s}[s]$.
Thus $\mathcal H_{a+1,b}=E(a)\mathcal H_{a,b}$ and
$\mathcal H_{a,b}=\mathcal H_{a,b+1}E(b)$.  Taking coefficients proves
the two shift formulas.
\end{proof}

The usual strict stuffle identity also holds without exclusions:
\begin{equation}\label{eht:stuffle-two}
 h_{a,b}(s)h_{a,b}(t)=
 \mathcal H_{a,b}(s,t)+\mathcal H_{a,b}(t,s)+h_{a,b}(s+t).
\end{equation}
More generally $\mathcal H_{a,b}$ is a character of the quasi-shuffle
algebra whose merged-letter operation is $[s]\circ[t]=[s+t]$.  Indeed $Z_a$ and
$Z_b$ are characters in the initial domain by splitting equality and
inequality cases in a product of sums.  This remains true meromorphically.
Characters into a commutative algebra are closed under convolution and
convolution inversion in the connected quasi-shuffle Hopf algebra
\cite{eht:Hoffman}; \eqref{eht:quotient} therefore gives the assertion.
Its coefficients are entire by Theorem~\ref{eht:main}.

Inverting the first shift formula gives the useful finite identity
\begin{equation}\label{eht:inverse-shift}
 \boxed{\mathcal H_{a,b}(s_1,\ldots,s_r)=
  \sum_{j=0}^r(-1)^j a^{-S_j}
                  \mathcal H_{a+1,b}(s_{j+1},\ldots,s_r),}
 \qquad S_0=0.
\end{equation}
It continues the upper endpoint to the principal slit plane
$\mathbb C\setminus(-\infty,0]$, by shifting it finitely many times into
the right half-plane.  Near $a=0$ the right side is a finite sum of
$a^{-S_j}$ times functions holomorphic in $a$.  Taylor expansion of those
functions gives a \emph{convergent} local expansion.  Thus the exact
branch exponents are prefix sums, and every fixed spectral derivative
adds only finitely many powers of $\log a$.

Continuation near $a=-N$, on a specified local branch, follows from
\[
 \mathcal H_{a,b}=E(a)^{-1}E(a+1)^{-1}\cdots E(a+N)^{-1}
                         \mathcal H_{a+N+1,b}.
\]
At fixed depth this is finite coefficient algebra.  Only the factor
$E(a+N)^{-1}$ is locally singular, so the branch exponents are sums over
consecutive subwords.  No essential singularity is introduced.  The
lower endpoint has the simpler local form \eqref{eht:lower-shift}: near
$b=0$ it has one possible singular factor $b^{-s_r}$.

\subsection{Gamma closure of the symmetric integer-index sector}

The symmetric-sum theorem for a stuffle character gives
\begin{equation}\label{eht:symmetric}
 \sum_{\sigma\in S_r}\mathcal H_{a,b}(s_{\sigma(1)},\ldots,s_{\sigma(r)})
 =\sum_{\pi}(-1)^{r-|\pi|}
       \prod_{B\in\pi}(|B|-1)!\,
       \prod_{B\in\pi}h_{a,b}\left(\sum_{j\in B}s_j\right),
\end{equation}
where $\pi$ runs over set partitions of $\{1,\ldots,r\}$.
This is the classical symmetric-sum mechanism applied to the entire
character.  Repeated labels are counted with their permutation
multiplicities; no independence of zeta values is assumed.

In particular, as a convergent germ at $u=0$,
\begin{equation}\label{eht:ones-Gamma}
 \boxed{\sum_{r\ge0}\mathcal H_{a,b}(\{1\}^r)u^r
       =\frac{\Gamma(a+u)\Gamma(b)}{\Gamma(a)\Gamma(b+u)}.}
\end{equation}
Its first coefficients, with $h_j=h_{a,b}(j)$, are
\[
 \mathcal H_{a,b}(1,1)=\frac{h_1^2-h_2}{2},\qquad
 \mathcal H_{a,b}(1,1,1)=\frac{h_1^3-3h_1h_2+2h_3}{6},
 \qquad h_1=\psi(a)-\psi(b).
\]

There is a simultaneous finite Gamma formula for every finite alphabet
of positive integer exponents.  Let $p_1,\ldots,p_d\ge1$, let
$M=\max_jp_j$, and let $\rho_1,\ldots,\rho_M$ be the roots, with
multiplicity, of
\[
 x^M+\sum_{j=1}^du_jx^{M-p_j}.
\]
Then
\begin{equation}\label{eht:alphabet-Gamma}
 \boxed{\sum_{r\ge0}\sum_{j_1,\ldots,j_r=1}^d
  \mathcal H_{a,b}(p_{j_1},\ldots,p_{j_r})u_{j_1}\cdots u_{j_r}
  =\prod_{\ell=1}^M
       \frac{\Gamma(a-\rho_\ell)\Gamma(b)}
            {\Gamma(a)\Gamma(b-\rho_\ell)}.}
\end{equation}
The right side is symmetric in the roots and therefore a holomorphic
germ in the $u_j$, even at colliding roots.  It evaluates sums over
permutations of the prescribed indices; it does not identify individual
ordered words with each other.

To prove \eqref{eht:alphabet-Gamma}, abelianize the quasi-shuffle series.
Its logarithm is the formal difference of
$\sum_n\log(1+\sum_ju_j(n+b)^{-p_j})$ and the corresponding $a$ series.
This statement first follows when all independent indices have real part
greater than one; the symmetric-sum identity and the entire differences
$h_{a,b}$ continue it to the integer alphabet.  Each coefficient of the
difference converges, including the linear exponent-one coefficient.
Factoring the polynomial gives the logarithm
\[
 -\sum_{m\ge1}\frac{h_{a,b}(m)}m
                     \sum_{\ell=1}^M\rho_\ell^m.
\]
The Taylor expansion of
$\log\Gamma(a-\rho)-\log\Gamma(a)
 -\log\Gamma(b-\rho)+\log\Gamma(b)$ is exactly this expression.
This proves the identity and also supplies a branch-independent
coefficient algorithm.

\subsection{Independent logarithmic jets and the ordering coordinate}

Put $\gamma_m(x)$ for the generalized Stieltjes constants in
\[
 \zeta(1+u,x)=\frac1u+
       \sum_{m\ge0}\frac{(-1)^m\gamma_m(x)}{m!}u^m.
\]
Then
\begin{equation}\label{eht:h-jets}
 h_{a,b}^{(m)}(1)=(-1)^m\{\gamma_m(b)-\gamma_m(a)\}.
\end{equation}
Differentiating \eqref{eht:stuffle-two} on the diagonal gives the exact
first-jet identity
\begin{align}\label{eht:sum-jets}
 (\partial_1+\partial_2)\mathcal H_{a,b}(1,1)
  ={}&\{\psi(a)-\psi(b)\}\{\gamma_1(a)-\gamma_1(b)\}\notag\\
    &-\zeta'(2,b)+\zeta'(2,a).
\end{align}

The incoming ordered-resonance paper uses the coefficient $\eta(a)$ in
\[
 Z_a(1+u,1)=u^{-2}+\gamma_0(a)u^{-1}
       +\frac{\gamma_0(a)^2-\zeta(2,a)}2+\eta(a)u+O(u^2).
\]
Substitution in \eqref{eht:depthtwo} now gives
\begin{align}
 \partial_1\mathcal H_{a,b}(1,1)
   &=\eta(b)-\eta(a)+\gamma_1(a)\{\psi(a)-\psi(b)\},
       \label{eht:first-eta}\\
 \partial_2\mathcal H_{a,b}(1,1)
   &=\eta(a)-\eta(b)-\gamma_1(b)\{\psi(a)-\psi(b)\}
                     -\zeta'(2,b)+\zeta'(2,a).
       \label{eht:second-eta}
\end{align}
These formulas explain the scope of the cancellation theorem precisely.
It removes the moving poles and makes the separate logarithmic
decorations well defined; it does not remove the nonsymmetric arithmetic
coordinate $\eta$.

\subsection{Finite reductions with a nonpositive inner exponent}

One independent-index subfamily does have complete depth-one closure.
For $m\ge0$ and arbitrary complex $s$,
\begin{equation}\label{eht:negative-inner}
 \boxed{\mathcal H_{a,b}(s,-m)=\frac1{m+1}
 \left\{\sum_{j=0}^{m+1}\binom{m+1}{j}B_{m+1-j}\,
                       h_{a,b}(s-j)
          -B_{m+1}(b)h_{a,b}(s)\right\}.}
\end{equation}
All apparent singular values in the separate Hurwitz functions are
removable within the $h$ differences.  In particular,
\begin{align}
 \mathcal H_{a,b}(s,0)&=h_{a,b}(s-1)-b\,h_{a,b}(s),
        \label{eht:zero-inner}\\
 \mathcal H_{a,b}(0,s)&=(a-1)h_{a,b}(s)-h_{a,b}(s-1),\notag\\
 \mathcal H_{a,b}(s,-1)&=\frac12\{h_{a,b}(s-2)-h_{a,b}(s-1)
                                -b(b-1)h_{a,b}(s)\}.\notag
\end{align}

For a direct proof, first take $\Re s>m+2$ and use the ordinary finite
power-sum identity
\[
 \sum_{k=0}^{n-1}(k+a)^m
       =\frac{B_{m+1}(n+a)-B_{m+1}(a)}{m+1}
\]
in $Z_a(s,-m)$.  Expand $B_{m+1}(n+a)$ in powers of $n+a$.
Insert the resulting finite Hurwitz formula, and its $b$ version, into
\eqref{eht:depthtwo}.  Since
$h_{a,b}(-m)=\{B_{m+1}(a)-B_{m+1}(b)\}/(m+1)$, the terms involving
$B_{m+1}(a)$ cancel.  This yields \eqref{eht:negative-inner}.  Entire
continuation gives every $s$.  The second formula in the displayed list
also follows from the stuffle identity.

As explicit Gamma--Stieltjes specializations at $b=1$,
\begin{align}
 \mathcal H_{a,1}(1,0)&=a-1-\psi(a)-\gamma,\notag\\
 \partial_1\mathcal H_{a,1}(1,0)
   &=-\log\Gamma(a)-\gamma_1(a)+\gamma_1,\label{eht:Gamma-Stieltjes}\\
 \partial_2\mathcal H_{a,1}(0,1)
   &=(a-1)\{\gamma_1(a)-\gamma_1\}+\log\Gamma(a).\notag
\end{align}
The derivatives are taken before the indicated index specialization.
The identity $h_{a,1}'(0)=-\log\Gamma(a)$ fixes the normalization.

These reductions also have normalized antiderivatives.  For a path in
the right half-plane from $b$ to $a$, define
\begin{equation}\label{eht:primitive}
 P_{a,b}(s)=\int_b^a h_{x,b}(s)\,dx
    =(a-b)\zeta(s,b)
     +\frac{\zeta(s-1,a)-\zeta(s-1,b)}{s-1}.
\end{equation}
The combined right side is entire in $s$.  At $s=1$ it equals
\[
 P_{a,b}(1)=\log\Gamma(a)-\log\Gamma(b)-(a-b)\psi(b).
\]
Replacing each $h_{a,b}(s-j)$ in \eqref{eht:negative-inner} by
$P_{a,b}(s-j)$ gives the uniquely normalized integral
$\int_b^a\mathcal H_{x,b}(s,-m)\,dx$.  Every spectral derivative
commutes with this identity.  Thus this independent-index family has a
finite calculus in Hurwitz-zeta derivatives, Gamma functions, and
generalized Stieltjes constants, with its additive constant fixed.

\subsection{Polylogarithmic dictionary}

For $|z|<1$, the interpolation theorem gives
\begin{equation}\label{eht:polylog}
 \sum_{N\ge0}\mathcal H_{N+1,1}(s_1,\ldots,s_r)z^N
  =\frac{\operatorname{Li}_{s_1,\ldots,s_r}(z)}{1-z},
\end{equation}
where the strict multiple polylogarithm has numerator $z^{n_1}$.
This is an ordinary normally convergent identity, locally in all
independent exponents.  Each spectral derivative inserts the expected
logarithmic factors in the harmonic sums, and the normalized Euler
primitive of the numerator raises its first index by one.  These
classical correspondences now attach to an entire interpolation in the
finite harmonic endpoint.

This correspondence also explains the compatibility with the positive-word
shuffle calculus in Section~\ref{mixed:section}: finite harmonic endpoint
coefficients and nested polylogarithmic words encode the same strict
ordering. The independent-order kernel and nonsymmetric-jet problems
are collected in Section~\ref{research:section}.
